\documentclass[11pt]{article}
\usepackage[margin=.85in,headheight=15pt]{geometry}
\usepackage[T1]{fontenc}
\usepackage{lmodern,microtype,mathtools,amssymb,booktabs,xcolor}
\usepackage{hyperref,xurl,fancyhdr,needspace}
\definecolor{ink}{HTML}{17324D}
\definecolor{teal}{HTML}{167D82}
\definecolor{muted}{HTML}{5D6B78}
\color{ink}
\hypersetup{colorlinks=true,urlcolor=teal,linkcolor=teal,
pdftitle={Presenter notes for two continued-fraction lectures},pdfauthor={Igor Rivin}}
\setlength{\parindent}{0pt}
\setlength{\parskip}{5pt}
\pagestyle{fancy}
\fancyhf{}
\fancyhead[L]{\small Continued fractions: presenter notes}
\fancyhead[R]{\small Two 50-minute lectures}
\fancyfoot[C]{\thepage}
\newcommand{\Q}{\mathbb Q}
\newcommand{\Z}{\mathbb Z}
\newcommand{\R}{\mathbb R}
\newcommand{\floor}[1]{\left\lfloor#1\right\rfloor}
\newcommand{\mat}[4]{\begin{pmatrix}#1&#2\\#3&#4\end{pmatrix}}
\newcounter{slide}
\newcounter{elapsed}
\newcounter{backup}
\newcommand{\teachslide}[4]{%
  \stepcounter{slide}\Needspace{9\baselineskip}
  \subsection*{\theslide.\ #1}
  {\small\color{teal}\textbf{Minutes \theelapsed\ to \number\numexpr\value{elapsed}+#2\relax}\par}
  \addtocounter{elapsed}{#2}
  #4\par}
\let\teachcover\teachslide
\newcommand{\backups}{%
  \typeout{LECTURE DURATION: \theelapsed\space minutes; CORE SLIDES: \theslide}
  \ifnum\value{elapsed}=50\else\errmessage{Lecture timing is not 50 minutes}\fi
  \setcounter{backup}{0}
  \Needspace{10\baselineskip}\section*{Optional slides after the 50-minute stopping point}}
\newcommand{\backupslide}[3]{%
  \stepcounter{backup}\Needspace{9\baselineskip}
  \subsection*{Backup \thebackup.\ #1}
  #3\par}
\begin{document}
{\huge\bfseries Continued fractions\par}
{\Large Presenter notes for two 50-minute lectures\par}
\vspace{10pt}
The slides follow Chapter 3 of \emph{Patterns, Proofs, and Primes}.
Each lecture has its own 50-minute route. The suggested times include
student calculation and discussion. Answer slides follow question slides,
so advancing the PDF provides a simple reveal without animation software.
The optional slides come after the exit question.

\textbf{Before the first lecture.}
Students should know how to multiply and invert a $2\times2$ matrix,
how a simple continued fraction is written, and what its convergents are.
Lecture 1 reviews the transformation rules.
For Lecture 2, recall that the convergents tend to the original
irrational number and their denominators are eventually unbounded.
No calculus or abstract algebra is needed for these proofs.

\textbf{Using the files.}
Share only the slide PDF window in Zoom and keep this notes document
separate. If you have one screen, practice switching to the notes
without changing the shared window. The editable LaTeX sources are
included in the source bundle. No live web services or fonts need to
download during the lectures.

\textbf{Pacing.}
Lecture 1 reaches the prefix-conjugation identity at minute 29 and
finishes that argument by minute 35.
Lecture 2 completes Lagrange's converse at minute 24, begins the
Galois discussion at minute 30, and completes its proof at minute 41.
If discussion runs long, shorten the final comparison examples.
Keep the student calculation breaks whenever possible.

\section*{Lecture 1: Continued fractions and matrix fixed points}
\teachcover{Continued fractions\\and matrix fixed points}{1}{
One repeating block, one quadratic equation
\[
  x=\frac{ax+b}{cx+d}
\]
}{
Open with the question: why should repeating continued fractions know
anything about quadratic equations? Students have already met linear
fractional transformations. Treat the first few slides as a review,
not a new chapter of prerequisites. Today's destination is the
prefix-conjugation argument. Source: course book, Sections 3.2 and 3.7.
}

\teachslide{A repeating pattern}{2}{
\[
 \sqrt2=[1;2,2,2,\ldots]=[1;\overline2]
\]
\[
 1+\sqrt2=[2;2,2,2,\ldots]=[\overline2]
\]
\spacer
The first expansion is \key{eventually periodic}.

The second is \key{purely periodic}.
}{
Ask students which digit belongs to the repeating block in each
expression. Explain that pure periodicity includes the integer part.
A bar over the block means to repeat it indefinitely. The original
number and its complete quotient can have different kinds of
periodicity. Recall that rational simple continued fractions terminate.
In these lectures, a finite rational expansion does not count as
eventually periodic. Source: course book, Section 3.7.
}

\teachslide{The two periodicity theorems}{2}{
A \key{quadratic irrational} is an irrational real root of an integer quadratic.
\spacer
\textbf{Lagrange}
\[
 \text{eventually periodic}
 \quad\Longleftrightarrow\quad
 \text{quadratic irrational}
\]
\spacer
\textbf{Galois}

A quadratic irrational $x$ is purely periodic exactly when
\[
 \key{x>1\qquad\text{and}\qquad -1<x'<0.}
\]
Here $x'$ is the other root of its quadratic equation.
}{
State the destination, without trying to prove both results now.
Use $\sqrt2$ as the simplest example of a quadratic irrational.
Today proves the forward direction of Lagrange. Lecture 2 proves the
converse and both directions of Galois. Clarify the names gently:
Legendre's continued-fraction theorem concerns exceptionally good
rational approximations. Lagrange is the eventual-periodicity theorem.
The theorem statements use simple continued fractions. Source:
course book, Sections 3.7--3.9. External cross-check: Corless,
\href{https://www.cecm.sfu.ca/organics/papers/corless/confrac/html/node6-an.shtml}{Periodic and Fixed Points of the Gauss Map}.
}

\teachslide{A M\"obius transformation}{2}{
\[
 M=\mat abcd,\qquad
 \key{T_M(t)=\frac{at+b}{ct+d}}
\]
\[
 a,b,c,d\in\Q,\qquad ad-bc\ne0
\]
\spacer
We use the formula where $ct+d\ne0$.

The determinant condition makes the transformation invertible.
}{
Use the notation from your preceding lecture, while making clear that
$T_M(t)$ is a fractional action rather than ordinary matrix
multiplication by a scalar. The denominator condition is essential.
With rational coefficients, any finite pole is rational, so an
irrational input never hits a pole. Infinity is a convenient optional
convention, not a prerequisite for today's proofs. Source: book 3.2.
}

\teachslide{Why a matrix acts on a number}{2}{
\[
 \mat abcd\begin{pmatrix}t\\1\end{pmatrix}
   =\begin{pmatrix}at+b\\ct+d\end{pmatrix}
\]
\spacer
Take the ratio of the two output coordinates.
\[
 t\ \longmapsto\ \frac{at+b}{ct+d}
\]
\spacer
Multiplying every entry of $M$ by the same nonzero number
leaves $T_M$ unchanged.
}{
This is the projective viewpoint without projective-space terminology.
Emphasize the ratio. The columns are ordinary vectors, and the
fractional transformation records their coordinate ratio.
Ask why $M$ and $-M$ produce the same map. The identity matrix and
all its nonzero scalar multiples therefore represent the identity
transformation. We will need that observation when excluding
degenerate fixed-point equations. Source: book 3.2 and Appendix A.2.
}

\teachslide{Composition}{3}{
\[
 \key{T_M\circ T_N=T_{MN}}
\]
\spacer
The rightmost matrix acts first.
\[
 M=\mat1101,\qquad N=\mat0110
\]
\[
 T_M(T_N(t))=1+\frac1t,\qquad
 MN=\mat1110
\]
}{
Here $M$ is translation by 1 and $N$ is inversion. Compute the
product with the class, then compare it with the nested function.
If students already remember the general substitution proof,
do not repeat its four coefficients. The first backup slide
contains that calculation if requested. Ask what happens in the
opposite order: $NM=\mat0111$ acts by $1/(t+1)$.
This makes the order of multiplication memorable. Source: book 3.2.
}

\teachslide{An inverse, by solving for $t$}{2}{
\[
 y=\frac{at+b}{ct+d}
 \qquad\Longrightarrow\qquad
 \key{t=\frac{dy-b}{a-cy}}
\]
\spacer
\[
 M^{-1}=\frac1{ad-bc}\mat d{-b}{-c}a
\]
\[
 (T_M)^{-1}=T_{M^{-1}}
\]
The scalar factor disappears when we take a coordinate ratio.
}{
Let students supply the rearrangement
$(a-cy)t=dy-b$. An important corollary is that an invertible
rational M\"obius transformation preserves irrationality:
if its output were rational, its inverse would make the input
rational. There is no hidden pole in this argument.
For $y=T_M(t)$, $a-cy=(ad-bc)/(ct+d)\ne0$.
We use this fact later to rule out a merely linear equation.
Source: book 3.2 and the proof in 3.7.
}

\teachslide{One continued-fraction digit}{2}{
\[
 \key{D(a)=\mat a110}
\]
\[
 T_{D(a)}(t)=a+\frac1t
\]
\spacer
The matrix inserts one digit before the remaining tail.
\[
 \det D(a)=-1,\qquad D(a)^{-1}=\mat011{-a}
\]
}{
The tail $t$ here is a whole remaining continued fraction.
The letter $D$ in $D(a)$ means a digit matrix. In Lecture 2,
the plain letter $D$ will instead denote a discriminant under
a square root; point out the contextual distinction if needed.
All finite digit products have determinant $\pm1$ and integer
inverses. This is the exact sense in which we can undo a prefix.
Source: book 3.2--3.3.
}

\teachslide{A block of digits}{2}{
\[
 [a_0;a_1,\ldots,a_{r-1},t]
 =T_{D(a_0)D(a_1)\cdots D(a_{r-1})}(t)
\]
\spacer
For the block $1,2$,
\[
 D(1)D(2)=\mat3121
\]
\[
 1+\frac1{2+1/t}=\frac{3t+1}{2t+1}.
\]
}{
Translate the displayed continued fraction back into nested
functions. All digits after the integer part are positive.
The last entry $t$ is the tail, so it need not be an integer.
The example fixes the period matrix that will run through the
rest of the lecture. Have students keep $\mat3121$ on their
paper. Source: book 3.2 and the worked example in 3.7.
}

\teachslide{A repeating tail}{2}{
\[
 x=[\overline{b_0,\ldots,b_{\ell-1}}],
 \qquad b_i\ge1
\]
\[
 N=D(b_0)\cdots D(b_{\ell-1})
   =\mat ABC{D_0}
\]
\spacer
Removing one whole period leaves the same tail.
\[
 \key{x=T_N(x)=\frac{Ax+B}{Cx+D_0}}
\]
Here $B,C>0$ and $\det N=(-1)^\ell$.
}{
This is the central idea. Pause before showing the fixed-point
equation and ask what remains after removing a complete block.
An infinite simple continued fraction has a well-defined
irrational limit; this follows from the convergent bounds
already developed in class. A backup gives the short
irrationality argument. Positivity of the period digits makes
$B,C$ strictly positive, even for a one-digit period.
Source: book 3.7.
}

\teachslide{A fixed point gives an equation}{2}{
\[
 x=\frac{Ax+B}{Cx+D_0}
\]
\[
 Cx^2+D_0x=Ax+B
\]
\[
 \key{Cx^2+(D_0-A)x-B=0}
\]
\spacer
Since $C>0$, this is a genuine quadratic equation.

Thus a purely periodic simple continued fraction is quadratic.
}{
Keep this algebraic step brief.
The nontrivial part will be understanding exactly which
quadratics repeat. At this point the number is irrational
because the continued fraction is infinite, so it is a
quadratic irrational, not merely some root of a reducible
quadratic. Contrast the period matrix with the identity:
$t=T_I(t)$ only yields $0=0$. Source: book 3.7.
}

\teachslide{Your turn: the block $1,2$}{3}{
\[
 x=[\overline{1,2}],\qquad N=\mat3121
\]
\spacer
Find a quadratic equation for $x$.

Find both roots.

Which root can be the continued fraction?
\spacer
\key{Take one minute before we compare answers.}
}{
Give students a real minute of silence, then ask for a
quadratic before asking for its roots. Expected equation:
$2x^2-2x-1=0$. Expected roots: $(1\pm\sqrt3)/2$.
The continued fraction exceeds 1, so it is the positive root.
If someone takes the negative root, ask which interval contains
$1+1/(2+1/(\cdots))$. Advance only after discussion.
This is book Exercise 15, expanded into a classroom prompt.
}

\teachslide{The periodic root}{2}{
\[
 2x^2-2x-1=0
\]
\[
 \key{x=[\overline{1,2}]=\frac{1+\sqrt3}{2}}
\]
\spacer
The other root is
\[
 x'=\frac{1-\sqrt3}{2},\qquad -1<x'<0.
\]
}{
The prime denotes the other algebraic root, not a derivative.
Do not yet explain the significance of the interval $(-1,0)$.
Point out the pattern and promise that Lecture 2 will
explain it exactly. One can check the value by substituting
it into $(3x+1)/(2x+1)$. Source: book 3.7.
}

\teachslide{A finite prefix}{2}{
\[
 x=T_N(x),\qquad y=T_M(x)
\]
\spacer
The original number $y$ may have extra digits before the period.
\[
 y=[1;\overline{1,2}]=1+\frac1x,
 \qquad M=D(1).
\]
\spacer
Which transformation fixes \key{$y$}?
}{
Emphasize the distinction between a periodic tail and the original
number. We know a quadratic equation for the tail
$x$, but must produce one for the number $y$ with its prefix.
Let students suggest substitution or inversion.
Direct substitution works, but conjugation packages it into
one reusable identity. Source: book 3.7, equation (3.3).
}

\teachslide{The prefix conjugation identity}{3}{
\[
 \begin{aligned}
 y&=T_M(x)\\[3pt]
  &=T_M\bigl(T_N(x)\bigr)\\[3pt]
  &=T_M\bigl(T_N(T_{M^{-1}}(y))\bigr).
 \end{aligned}
\]
\spacer
\[
 \key{y=T_{MNM^{-1}}(y)}
\]
Attaching a prefix conjugates the transformation that fixes the tail.
}{
Read the composition from right to left: remove the prefix,
apply the period to its fixed tail, and put the prefix back.
Spend the full three minutes on why each equality holds.
For students comfortable with function notation, write
$MNM^{-1}$ directly while saying that the matrices act
fractionally. An empty prefix means $M=I$.
Source: book 3.7, equation (3.3).
}

\teachslide{Why the new equation is genuinely quadratic}{3}{
Put $K=MNM^{-1}$.

\spacer
$K$ has integer entries and determinant $\pm1$.

$K$ is nonscalar, since conjugating a scalar matrix gives a scalar matrix.

$y$ is irrational, since $T_{M^{-1}}(y)=x$.
\spacer
For $K=\mat abcd$,
\[
 \key{cy^2+(d-a)y-b=0.}
\]
This polynomial is nonzero and cannot be merely linear.
}{
This is the rigor checkpoint. If $K=\lambda I$, then
$N=M^{-1}KM=\lambda I$, contradicting the positive
lower-left entry of the period matrix. Thus the three
coefficients $c,d-a,-b$ are not all zero.
If $c=0$, a nonzero rational linear polynomial would vanish
at the irrational $y$, which is impossible.
Rationality of a M\"obius image would also imply rationality
of its inverse image. The backup fixed-point lemma gives
this argument in isolation. Source: book 3.7.
}

\teachslide{The equation for $\sqrt3$}{3}{
\[
 M=\mat1110,\qquad
 N=\mat3121,\qquad
 M^{-1}=\mat011{-1}
\]
\[
 MNM^{-1}=\mat2312
\]
\[
 y=\frac{2y+3}{y+2}
 \qquad\Longrightarrow\qquad y^2=3
\]
\[
 \key{[1;\overline{1,2}]=\sqrt3}
\]
}{
Have students check at least one row of the matrix product.
For an intermediate check, $MN=\mat5231$.
The positive sign follows from $y=1+1/x>0$.
This example both verifies the familiar continued fraction
of $\sqrt3$ and shows what conjugation actually does to
the quadratic equation. Source: book 3.7 and Exercise 21.
}

\teachslide{Your turn: a different prefix}{3}{
Keep $x=[\overline{1,2}]$, but use $M=D(2)$.
\[
 y=[2;\overline{1,2}]=2+\frac1x
\]
\spacer
Compute $MNM^{-1}$ and identify $y$.

Can you write its digits with a bar starting at the first digit?
\spacer
\key{A written prefix need not be a genuine preperiod.}
}{
Allow about 90 seconds for the calculation.
Students who know $1/x=\sqrt3-1$ can immediately identify
$y=1+\sqrt3$, but ask them to check the conjugated matrix too.
The result is $\mat3211$, giving $y^2-2y-2=0$.
The digit stream is $2,1,2,1,\ldots$, so its pure period
is $[\overline{2,1}]$. Source: book Exercise 21.
}

\teachslide{A prefix that disappears}{2}{
\[
 D(2)ND(2)^{-1}=\mat3211
\]
\[
 y^2-2y-2=0,\qquad y=1+\sqrt3
\]
\spacer
\[
 \key{[2;\overline{1,2}]=[\overline{2,1}]}
\]
\spacer
The digits repeat from the beginning after all.
}{
Distinguish a particular written decomposition into prefix
and period from the shortest possible preperiod of the
actual digit sequence. Here the first digit already belongs
to a cycle of length two. This is an important example to
bring back after Galois's criterion, since the conjugate
$1-\sqrt3$ lies in $(-1,0)$. Source: book 3.9.
}

\teachslide{Quadratic conjugates}{2}{
\[
 x=u+v\sqrt D,\qquad
 \key{x'=u-v\sqrt D}\qquad (u,v\in\Q,\ v\ne0)
\]
\spacer
Replacing $\sqrt D$ by $-\sqrt D$ respects rational arithmetic.
\[
 \bigl(T_M(x)\bigr)'=T_M(x')
\]
\spacer
For example,
\[
 (\sqrt2)'=-\sqrt2,\qquad
 (1+\sqrt2)'=1-\sqrt2.
\]
}{
Assume $D$ is a positive nonsquare integer.
Clarify again that the prime is algebraic conjugation,
not differentiation. All we need is that changing the
sign of $\sqrt D$ commutes with adding, multiplying,
and taking reciprocals when the denominators are nonzero.
The complete quotients are rational transformations of
the original surd, so their conjugates follow the same
algebra. Source: book 3.7--3.8.
}

\teachslide{Reduced surds}{3}{
A quadratic irrational is \key{reduced} when
\[
 \key{x>1\qquad\text{and}\qquad -1<x'<0.}
\]
\spacer
\begin{center}
\renewcommand{\arraystretch}{1.25}
\begin{tabular}{lll}
\toprule
$x$ & $x'$ & Reduced?\\
\midrule
$\sqrt2$ & $-\sqrt2$ & No\\
$1+\sqrt2$ & $1-\sqrt2$ & Yes\\
$\sqrt2/2$ & $-\sqrt2/2$ & No\\
\bottomrule
\end{tabular}
\end{center}
}{
Ask which inequality fails in each of the two negative
examples. For $\sqrt2$ it is the lower bound on the
conjugate. For $\sqrt2/2$ it is $x>1$.
Galois says that this definition is exactly the right
one for pure periodicity. Today students need the definition
and examples; Lecture 2 explains the mechanism.
Source: book 3.8--3.9.
}

\teachslide{What we have proved}{2}{
\[
 \key{\text{eventually periodic}\quad\Longrightarrow\quad
 \text{quadratic irrational}}
\]
\spacer
A periodic tail satisfies a fixed-point equation.

A finite prefix changes $N$ to $MNM^{-1}$.
\spacer
Next lecture: why every quadratic irrational repeats,
and why reduced surds repeat immediately.
\spacer
\textbf{Exit question:} Why does $t=T_I(t)$ prove nothing?
}{
Expected exit answer: the identity transformation fixes
every real number, so the resulting polynomial is $0=0$.
We needed a nonscalar transformation, supplied by the
positive lower-left entry of the period matrix.
Stop here at 50 minutes. The remaining slides are backups,
not part of the planned lecture. Suggested preparation:
review complete quotients and the fact that convergents
$p_k/q_k$ tend to the original irrational.
}

\backups

\backupslide{Composition in coordinates}{
\[
 M=\mat abcd,\qquad N=\mat efgh
\]
\[
 \begin{aligned}
 T_M(T_N(t))
 &=\frac{a(et+f)+b(gt+h)}{c(et+f)+d(gt+h)}\\[8pt]
 &=\frac{(ae+bg)t+(af+bh)}{(ce+dg)t+(cf+dh)}.
 \end{aligned}
\]
\spacer
These are exactly the four entries of $MN$.
}{
Use only if students request the general composition
calculation. The equation holds wherever the displayed
finite formulas are defined, and globally if infinity
is included in the projective convention. Source: book 3.2.
}

\backupslide{The nonscalar fixed-point lemma}{
Let $M=\mat abcd$ be an invertible rational matrix.
\[
 t=T_M(t)\quad\Longrightarrow\quad
 ct^2+(d-a)t-b=0.
\]
If all coefficients vanish, then $M=aI$.

If $c=0$ and the polynomial is nonzero, its root is rational.
\spacer
Therefore an irrational fixed point of a nonscalar rational
matrix is a quadratic irrational.
}{
This separates the elementary lemma from any special
continued-fraction properties. Note that being invertible
does not by itself exclude the identity transformation.
One needs the nonscalar assumption as well as an
irrational fixed point. Source: book Lemma 3.12.
}

\backupslide{Why an infinite simple continued fraction is irrational}{
The convergents are in lowest terms, $q_k\to\infty$, and
\[
 \left|x-\frac{p_k}{q_k}\right|
 \le\frac1{q_kq_{k+1}}.
\]
If $x=p/q$ were rational, then eventually $p_k/q_k\ne p/q$,
and
\[
 \left|\frac pq-\frac{p_k}{q_k}\right|
 \ge\frac1{qq_k}.
\]
For $q_{k+1}>q$, these inequalities contradict each other.
}{
The alternating-convergent argument produces a limit
for any infinite positive digit sequence, not just
for the expansion of a number already known to be irrational.
For a rational candidate $p/q$ in lowest terms, increasing
denominators ensure that the convergents eventually
differ from it. The integer numerator of their difference
then has absolute value at least one. A weak upper bound
is enough here. Source: book 3.7.
}

\clearpage
\setcounter{slide}{0}\setcounter{elapsed}{0}
\section*{Lecture 2: Lagrange, Galois, and reduced surds}
\teachcover{Lagrange, Galois,\\and reduced surds}{1}{
Which continued fractions repeat?
\[
 x>1,\qquad -1<x'<0
\]
}{
Begin with a quick reminder of the end of Lecture 1:
a repeating tail is a fixed point, and attaching a prefix
conjugates its matrix. Today's new work proves the converse
and the exact pure-periodicity criterion. Students need the
complete-quotient algorithm and convergence of $p_k/q_k$,
but no algebraic number theory. Source: book 3.8--3.9.
}

\teachslide{Recall: the complete quotient}{2}{
\[
 x_0=x,\qquad a_n=\floor{x_n},\qquad
 \key{x_{n+1}=\frac1{x_n-a_n}}
\]
\spacer
The number $x_n=[a_n;a_{n+1},\ldots]$ is the remaining tail.
\[
 x=[a_0;a_1,\ldots,a_{n-1},x_n].
\]
For $n\ge1$,
\[
 x_n>1,\qquad a_n\ge1.
\]
}{
Rebuild the notation rather than assuming everyone remembers
the subscript convention. The integer part of the original
number can be zero or negative. Every later complete quotient
is greater than 1, since it is the reciprocal of a nonzero
fractional part. A repeated complete quotient makes all future
digits repeat because the algorithm is deterministic.
Source: book 3.5 and 3.8.
}

\teachslide{The target theorem}{2}{
\textbf{Lagrange: the direction still to prove}
\[
 \key{\text{quadratic irrational}\quad\Longrightarrow\quad
 \text{eventually periodic}}
\]
\spacer
\textbf{Galois}
\[
 \key{\text{purely periodic}\quad\Longleftrightarrow\quad
 x>1,\ -1<x'<0}
\]
The second statement concerns quadratic irrationals.
\spacer
We will use integer states and a reversible step.
}{
Recall that $x'$ means the other root of the rational
quadratic equation, or the value obtained by reversing
the sign of its square root. The reducedness conditions
are not a conclusion about arbitrary real numbers;
they are defined for quadratic irrationals.
The key logical distinction is eventual versus pure
periodicity. Source: book Theorems 3.13 and 3.17.
External cross-check: Corless, \href{https://www.cecm.sfu.ca/organics/papers/corless/confrac/html/node6-an.shtml}{Periodic and Fixed Points of the Gauss Map}.
}

\teachslide{Conjugates follow the same algebra}{2}{
\[
 x_{n+1}=\frac1{x_n-a_n}
 \qquad\Longrightarrow\qquad
 \key{x_{n+1}'=\frac1{x_n'-a_n}}
\]
\spacer
The digit is still \key{$a_n=\floor{x_n}$}.

We do not take the integer part of $x_n'$.
\spacer
If $x_n'<0$ and $a_n\ge1$, then
\[
 x_n'-a_n<-1,\qquad -1<x_{n+1}'<0.
\]
}{
This is a common source of confusion. Conjugation changes
the algebraic number but leaves the rational integer $a_n$
fixed. It does not commute with taking a floor.
The inequality uses the floor of the positive complete
quotient. Ask students why the reciprocal of a number
less than $-1$ lies in $(-1,0)$. Once a negative
conjugate appears, the next tail is reduced.
Source: book 3.8, Step 1.
}

\teachslide{Reduced tails stay reduced}{2}{
Suppose
\[
 x_n>1,\qquad -1<x_n'<0.
\]
Then
\[
 x_{n+1}>1,\qquad -1<x_{n+1}'<0.
\]
\spacer
\key{Once reduced, every later complete quotient is reduced.}
\spacer
We still have to show that an arbitrary quadratic irrational
eventually reaches this region.
}{
The first inequality follows from the fractional-part
algorithm. The second follows from the previous slide.
Do not skip the last sentence: forward invariance is not
the same as eventual entry. This logical separation also
prevents a circular proof of Lagrange's theorem.
Source: book 3.8, Step 1.
}

\teachslide{Every quadratic irrational reaches a reduced tail}{3}{
The inverse prefix matrix gives, for $n\ge2$,
\[
 \key{x_n'=-\frac{q_{n-2}x'-p_{n-2}}
                    {q_{n-1}x'-p_{n-1}}.}
\]
The convergents tend to $x$, and $x'\ne x$.

For large $n$, numerator and denominator have the same sign.
\[
 x_n'<0\quad\text{eventually}.
\]
One more step puts the conjugate in $(-1,0)$.
}{
Give the full sign argument aloud. Since $p_k/q_k\to x$
and $x'\ne x$, the difference $x'-p_k/q_k$ eventually
has the fixed sign of $x'-x$. Multiplication by positive
$q_k$ preserves that sign. The two expressions in the
ratio therefore have the same sign, and the minus sign
makes $x_n'$ negative. Lecture 1's inverse formula gives
the displayed identity. A backup derives it in detail.
This covers negative original numbers too.
Source: book 3.8, Step 1.
}

\teachslide{An integer description of a surd}{2}{
Choose a fixed positive nonsquare integer $D$.
\[
 \key{x_n=\frac{P_n+\sqrt D}{Q_n}},
 \qquad P_n,Q_n\in\Z,\quad Q_n\ne0.
\]
We can choose the initial representation so that
\[
 \key{Q_n\mid D-P_n^2}
\]
holds at every step.
\spacer
$D$ need not be squarefree. Initially, $Q_n$ may be negative.
}{
Explain what the state remembers: the fixed square root
never changes, so only two integers need to evolve.
For completeness, if $Ax^2+Bx+C=0$, choose the common
sign of the coefficients so that
$x=(-B+\sqrt{B^2-4AC})/(2A)$.
Then $D=B^2-4AC$, $P_0=-B$, $Q_0=2A$, and
$D-P_0^2=-4AC$ is divisible by $Q_0$.
Do not automatically insist that the initial denominator
be positive. The first backup shows this normalization.
Source: book 3.8, Step 2.
}

\teachslide{The exact recurrence}{3}{
Set $a_n=\floor{x_n}$ and
\[
 \key{P_{n+1}=a_nQ_n-P_n.}
\]
Then
\[
 \begin{aligned}
 \frac1{x_n-a_n}
 &=\frac{Q_n}{\sqrt D-P_{n+1}}\\[5pt]
 &=\frac{P_{n+1}+\sqrt D}{(D-P_{n+1}^2)/Q_n}.
 \end{aligned}
\]
Thus
\[
 \key{Q_{n+1}=\frac{D-P_{n+1}^2}{Q_n}.}
\]
}{
Rationalize the denominator explicitly.
The sign in $P_{n+1}=a_nQ_n-P_n$ is easy to reverse
accidentally, so check it against
$x_n-a_n=(\sqrt D-P_{n+1})/Q_n$.
After rationalization, the numerator has precisely the
same form as before. We have not yet proved the new
denominator is an integer; that is the next slide.
Source: book 3.8, equation (3.4).
}

\teachslide{Why the recurrence stays integral}{2}{
\[
 P_{n+1}^2-P_n^2
 =a_nQ_n(a_nQ_n-2P_n).
\]
So $Q_n$ divides both $D-P_n^2$ and $D-P_{n+1}^2$.
\[
 \key{Q_{n+1}=(D-P_{n+1}^2)/Q_n\in\Z.}
\]
Also,
\[
 D-P_{n+1}^2=Q_nQ_{n+1}.
\]
This supplies the divisibility condition for the next step.
}{
This version uses divisibility directly, so students need
not already know congruence notation. Since $D$ is a
nonsquare integer, $D-P_{n+1}^2$ is never zero.
Thus $Q_{n+1}$ is a nonzero integer. The final equality
says that $Q_{n+1}$ itself divides
$D-P_{n+1}^2$, completing the induction.
Source: book 3.8, Step 2.
}

\teachslide{The bounds from reducedness}{3}{
For a reduced state,
\[
 x=\frac{P+\sqrt D}{Q}>1,\qquad
 -1<x'=\frac{P-\sqrt D}{Q}<0.
\]
Subtract:
\[
 x-x'=\frac{2\sqrt D}{Q}>1
 \quad\Longrightarrow\quad \key{0<Q<2\sqrt D.}
\]
The conjugate inequalities give
\[
 \sqrt D-Q<P<\sqrt D
 \quad\Longrightarrow\quad \key{|P|<\sqrt D.}
\]
}{
This is the compactness argument entirely in elementary
integer language. The positive difference $x-x'$ first
forces $Q>0$, so it is safe to multiply the conjugate
inequalities by $Q$. The lower bound becomes
$P>\sqrt D-Q>-\sqrt D$ because $Q<2\sqrt D$.
No uniform bound on the period length is needed.
Source: book 3.8, equation (3.5).
}

\teachslide{Finitely many states}{2}{
For fixed $D$, only finitely many integer pairs satisfy
\[
 \key{|P|<\sqrt D,\qquad 0<Q<2\sqrt D.}
\]
Every sufficiently late state lies in this finite set.

Each state determines its digit and its next state.
\[
 (P_i,Q_i)=(P_j,Q_j)
 \quad\Longrightarrow\quad
 a_{i+k}=a_{j+k}\quad(k\ge0).
\]
\key{This proves the converse direction of Lagrange's theorem.}
}{
Name the pigeonhole principle. Repetition of a pair means
repetition of the complete quotient, and deterministic
updating reproduces the entire future, not just one digit.
We have now proved both directions of Lagrange using
fixed points, convergence, rationalization, and finite
integer bounds. We still have not removed an initial
segment for a reduced starting state. Source: book 3.8.
}

\teachslide{The cycle for $\sqrt3$}{2}{
\begin{center}
\renewcommand{\arraystretch}{1.4}
\begin{tabular}{ccccc}
\toprule
$n$ & $P_n$ & $Q_n$ & $x_n$ & $a_n$\\
\midrule
$0$ & $0$ & $1$ & $\sqrt3$ & $1$\\
$1$ & $1$ & $2$ & $(1+\sqrt3)/2$ & $1$\\
$2$ & $1$ & $1$ & $1+\sqrt3$ & $2$\\
$3$ & $1$ & $2$ & $(1+\sqrt3)/2$ & $1$\\
\bottomrule
\end{tabular}
\end{center}
\[
 \key{\sqrt3=[1;\overline{1,2}]}
\]
The initial state lies outside the reduced cycle.
}{
Do one transition explicitly:
from $(P,Q,a)=(0,1,1)$ we get $P_{\rm new}=1$,
$Q_{\rm new}=2$. From $(1,2,1)$ we get $(1,1)$.
The equality of states at $n=1$ and $n=3$ certifies
period two. The finite table illustrates the proof,
but the bounds explain why every quadratic irrational
must eventually produce such a repetition.
Source: book Example 3.16.
}

\teachslide{Your turn: $\sqrt6$}{3}{
Start with $D=6$ and $(P_0,Q_0)=(0,1)$.
\[
 P_{n+1}=a_nQ_n-P_n,\qquad
 Q_{n+1}=\frac{6-P_{n+1}^2}{Q_n}.
\]
\spacer
Find the first repeated state.

Read off the continued fraction of $\sqrt6$.
\spacer
\key{Record the triples $(P_n,Q_n,a_n)$.}
}{
Give about two minutes of working time and one minute
to compare. The first digit is 2. The states are
$(0,1)$, $(2,2)$, $(2,1)$, $(2,2)$, with digits
$2,2,4,2$. If a student gets a negative denominator
after the first step, check the sign of $P_{n+1}$.
The interval $2<\sqrt6<3$ is enough to determine all
digits in this example without a decimal calculator.
This is a new worked classroom example using book (3.4).
}

\teachslide{The period of $\sqrt6$}{1}{
\begin{center}
\renewcommand{\arraystretch}{1.35}
\begin{tabular}{cccc}
\toprule
$n$ & $P_n$ & $Q_n$ & $a_n$\\
\midrule
$0$ & $0$ & $1$ & $2$\\
$1$ & $2$ & $2$ & $2$\\
$2$ & $2$ & $1$ & $4$\\
$3$ & $2$ & $2$ & $2$\\
\bottomrule
\end{tabular}
\end{center}
\[
 \key{\sqrt6=[2;\overline{2,4}]}
\]
}{
Keep the answer reveal brief. The first two reduced
complete quotients are $(2+\sqrt6)/2$ and $2+\sqrt6$.
Their conjugates are $(2-\sqrt6)/2$ and $2-\sqrt6$,
both in $(-1,0)$. This example gives a quick independent
check that the recurrence and reducedness conditions
are working together.
}

\teachslide{Finiteness allows an initial segment}{2}{
A function on the finite set $\{0,1\}$ can satisfy
\[
 F(0)=1,\qquad F(1)=1.
\]
Starting at 0 gives
\[
 0,\ 1,\ 1,\ 1,\ldots
\]
\spacer
Finiteness guarantees eventual repetition.

\key{To prove pure periodicity, we need to recover predecessors.}
}{
This is the conceptual turning point of the second lecture.
The reduced states form a finite forward-invariant set,
but that alone would not prove that an orbit starts on
its cycle. The little finite-map example exhibits the
missing possibility. Galois's criterion succeeds because
the continued-fraction step has unique reduced
predecessors. Source: book 3.9.
}

\teachslide{Galois's criterion}{2}{
For a quadratic irrational $x$,
\[
 \key{x\text{ is purely periodic}}
\]
\[
 \Longleftrightarrow
\]
\[
 \key{x>1\qquad\text{and}\qquad -1<x'<0.}
\]
\spacer
We first prove that a reduced starting point has no preperiod.
}{
Recall the definition from Lecture 1. We already know
that every reduced orbit stays reduced and eventually
repeats. The next two slides strengthen that conclusion
to pure periodicity. Afterward we prove necessity.
Keep the logical direction explicit throughout.
Source: book Theorem 3.17.
}

\teachslide{A successor remembers its reduced predecessor}{3}{
Suppose $u$ is reduced, $a=\floor u$, and
\[
 z=\frac1{u-a}.
\]
Conjugating gives
\[
 -\frac1{z'}=a-u',\qquad
 a<-\frac1{z'}<a+1.
\]
Therefore
\[
 \key{a=\floor{-1/z'},\qquad u=a+\frac1z.}
\]
Both the previous digit and the reduced predecessor are unique.
}{
The strict inequalities use precisely $-1<u'<0$.
This is why the definition of reducedness is so well
chosen. Stress that the assertion concerns reduced
predecessors. Without that restriction, many different
integers could be placed before the same tail.
For a quick check, take $z=(1+\sqrt3)/2$.
Then $-1/z'=1+\sqrt3$, so $a=2$ and the predecessor
is $1+\sqrt3$. Source: book 3.9, equation (3.6).
}

\teachslide{Rewinding a repeated state}{2}{
Suppose all $x_n$ are reduced and $x_i=x_j$ with $i<j$.
\[
 \begin{aligned}
 x_i&=x_j\\
 x_{i-1}&=x_{j-1}\\
 &\ \vdots\\
 x_0&=x_{j-i}.
 \end{aligned}
\]
Every backward step uses the unique reduced predecessor.
\spacer
\key{The complete-quotient sequence repeats from its start.}
}{
If $i=0$, the conclusion already holds. Otherwise apply
the predecessor formula repeatedly until reaching index
zero. Equal quadratic numbers have equal conjugates,
so the recovered digit and predecessor agree.
Now deterministic forward updating gives pure periodicity.
The point is not that we ran the ordinary floor algorithm
backward on arbitrary reals, but that reduced surds
carry enough information in their conjugates.
Source: book 3.9.
}

\teachslide{Why a pure period must be reduced}{2}{
Suppose the pure period has length $\ell$.

Lecture 1 shows that $x$ is quadratic.

All sufficiently late complete quotients are reduced.
\[
 \key{x_{k\ell}=x_0\qquad(k=1,2,\ldots)}
\]
Choose $k$ large enough.
\spacer
\key{The starting value is itself one of those reduced tails.}
}{
This is the other implication of Galois.
There is no circular appeal to Galois in the eventual-entry
argument: it used convergence of the ordinary convergents
and the inverse prefix matrix.
Together with the preceding two slides, this completes
the characterization. Source: book 3.9.
}

\teachslide{A reducedness test}{3}{
Which of these have purely periodic continued fractions?
\[
 \sqrt2,\qquad
 1+\sqrt3,\qquad
 \frac{1+\sqrt5}{2},\qquad
 \frac{\sqrt2}{2}
\]
\spacer
Check \key{both} inequalities:
\[
 x>1,\qquad -1<x'<0.
\]
\spacer
Which of the four are eventually periodic?
}{
Allow about 90 seconds, then ask for reasons rather than
just yes/no answers. The middle two are reduced.
For $\sqrt2$ the conjugate is below $-1$.
For $\sqrt2/2$ the original number is below 1.
All four are quadratic irrationals, so all four are
eventually periodic by Lagrange. Do not reveal the
next slide until the discussion finishes.
Source: book 3.9 and Exercise 23.
}

\teachslide{The test, checked}{2}{
\begin{center}
\renewcommand{\arraystretch}{1.4}
\begin{tabular}{lll}
\toprule
$x$ & $x'$ & Purely periodic?\\
\midrule
$\sqrt2$ & $-\sqrt2$ & No\\
$1+\sqrt3$ & $1-\sqrt3$ & Yes\\
$(1+\sqrt5)/2$ & $(1-\sqrt5)/2$ & Yes\\
$\sqrt2/2$ & $-\sqrt2/2$ & No\\
\bottomrule
\end{tabular}
\end{center}
\spacer
\key{All four are eventually periodic.}
}{
For the golden ratio, the inequalities follow from
$2<\sqrt5<3$. Its period consists of the digit 1.
For $1+\sqrt3$, the period is $2,1$.
Ask a student to distinguish the roles of the two theorems
in the final line. Source: book 3.9.
}

\teachslide{A visible prefix can belong to the period}{1}{
\[
 \key{[2;\overline{1,2}]=1+\sqrt3=[\overline{2,1}]}
\]
\spacer
Its conjugate is $1-\sqrt3\in(-1,0)$.

Galois detects pure periodicity independently of how
we first write the digit string.
}{
Return to the example from Lecture 1.
The earlier calculation established the value.
Galois now explains the classification.
A written bar need not mark the earliest possible
start of repetition. Source: book 3.9.
}

\teachslide{Lagrange and Galois}{2}{
\textbf{Lagrange}
\[
 \text{eventually periodic}
 \quad\Longleftrightarrow\quad
 \text{quadratic irrational}
\]
Fixed points prove one direction. Bounded integer states prove the other.
\spacer
\textbf{Galois}
\[
 \text{purely periodic}
 \quad\Longleftrightarrow\quad
 x>1,\ -1<x'<0
\]
Unique reduced predecessors remove the initial segment.
}{
Keep the quadratic-irrational hypothesis on the Galois
line explicit orally. The two lectures together used
only matrix algebra, convergence of convergents,
rationalization, inequalities, and the pigeonhole
principle. Mention optional links to Pell's equation
and reversal of the period only if time remains.
Both appear in the backup slides. Source: book 3.7--3.10.
}

\teachslide{Exit question}{1}{
In the proof, where did we use
\[
 \key{\text{finiteness?}}
\]
Where did we use
\[
 \key{\text{reversibility?}}
\]
\spacer
What can go wrong if only the first property holds?
}{
Expected answer: finiteness forces some repeated state
and therefore eventual periodicity. Unique reduced
predecessors let us rewind that repetition to index zero,
which gives pure periodicity. Without reversibility,
an orbit can have a transient initial segment, as in
$0,1,1,\ldots$. This is the 50-minute stopping point.
Suggested follow-up: book Exercises 24--26, or derive
the period of $\sqrt{13}$ using the exact recurrence.
}

\backups

\backupslide{Choosing the initial integer state}{
Start with $Ax^2+Bx+C=0$, where $A,B,C\in\Z$ and $A\ne0$.

Change all coefficient signs if needed so that
\[
 x=\frac{-B+\sqrt D}{2A},\qquad D=B^2-4AC>0.
\]
Then choose
\[
 P_0=-B,\qquad Q_0=2A.
\]
\[
 D-P_0^2=-4AC=Q_0(-2C).
\]
Thus $Q_0\mid D-P_0^2$. Its sign need not be positive.
}{
The discriminant is nonsquare because $x$ is irrational.
If the desired root initially has the minus sign before
the square root, replacing $(A,B,C)$ by $(-A,-B,-C)$
puts that root in the displayed plus-sign form.
For example, $x=-\sqrt2$ can be represented with
$D=8$, $P_0=0$, $Q_0=-2$.
This explains why a proof for all real quadratic
irrationals must allow initial negative denominators.
Source: book 3.8, Step 2.
}

\backupslide{The inverse-prefix formula}{
\[
 x=\frac{p_{n-1}x_n+p_{n-2}}
          {q_{n-1}x_n+q_{n-2}}.
\]
Solving for $x_n$ gives
\[
 x_n=-\frac{q_{n-2}x-p_{n-2}}
               {q_{n-1}x-p_{n-1}}.
\]
Conjugating gives
\[
 \key{x_n'=-\frac{q_{n-2}x'-p_{n-2}}
                     {q_{n-1}x'-p_{n-1}}.}
\]
For large $k$, $q_kx'-p_k=q_k(x'-p_k/q_k)$ has a fixed sign.
}{
Use this if the sign argument in main slide 6 needs
more time. The denominators are nonzero because the
conjugate is irrational and the convergents are rational.
For $n\ge2$, the relevant denominators $q_{n-2}$ and
$q_{n-1}$ are positive. Convergence toward $x$, not
any assumption of periodicity, gives the common sign.
Source: book 3.8, Step 1.
}

\backupslide{The conjugate reverses the period}{
If $x=[\overline{b_0,\ldots,b_{\ell-1}}]$, then
\[
 \key{-\frac1{x'}=[\overline{b_{\ell-1},\ldots,b_0}].}
\]
Each digit matrix is symmetric, so reversing a product
transposes it.
\[
 N=\mat ABC{D_0}:
 \quad Cx^2+(D_0-A)x-B=0
\]
Substitute $x'=-1/z$:
\[
 Bz^2+(D_0-A)z-C=0,
\]
the fixed-point equation for $N^T$.
}{
Both $B$ and $C$ are positive. The transformed quadratic
therefore has one positive root and one negative root.
By Galois, $-1/x'>0$; the reversed periodic continued
fraction is positive too, so they are the same root.
For a concrete check,
$[\overline{1,2,3}]=(4+\sqrt{37})/7$, while the reversed
period is $(4+\sqrt{37})/3$.
Source: book Corollary 3.19 and Exercise 26.
}

\backupslide{The same matrix produces Pell solutions}{
Lecture 1 gave
\[
 U=\mat2312,\qquad T_U(\sqrt3)=\sqrt3.
\]
Ordinary matrix multiplication gives
\[
 U\begin{pmatrix}p\\q\end{pmatrix}
   =\begin{pmatrix}2p+3q\\p+2q\end{pmatrix}.
\]
\[
 (2p+3q)^2-3(p+2q)^2=p^2-3q^2.
\]
Starting from $(1,0)$ produces
\[
 (2,1),\ (7,4),\ (26,15),\ldots
\]
Each satisfies $p^2-3q^2=1$.
}{
This is an optional connection, not an additional theorem
needed for either periodicity proof. The determinant-one
matrix acts both as a M\"obius transformation fixing
$\sqrt3$ and as an ordinary linear transformation
preserving the quadratic form $p^2-3q^2$.
Avoid claiming that this calculation alone proves the
classification of all Pell solutions.
Source: book 3.10 and Exercise 18.
}

\clearpage
\section*{Sources and notation}
The primary source for these lectures is the expanded course book:
Chapter 3, especially Sections 3.2 and 3.7--3.10, together with
Appendix A. The printed book numbers its three central identities
as (3.3), prefix conjugation; (3.4), the quadratic-state recurrence;
and (3.6), the unique reduced predecessor.

The notation $T_M(t)$ means the fractional action of a matrix,
not ordinary multiplication by a scalar. The prime on a quadratic
irrational denotes its algebraic conjugate, not a derivative.
Every continued fraction in the lectures is simple, with positive
digits after its integer part.

Robert M. Corless, \emph{Continued Fractions and Chaos},
\href{https://www.cecm.sfu.ca/organics/papers/corless/confrac/html/node6-an.shtml}{Periodic and Fixed Points of the Gauss Map},
states the Lagrange and Galois characterizations.
Giovanni Panti, \href{https://arxiv.org/abs/0712.2996}{A general Lagrange Theorem}
(2007), gives a broader setting for periodicity results beyond
ordinary simple continued fractions. The lectures need none of
those generalizations.

For a presentation setup reference, Zoom documents
\href{https://support.zoom.com/hc/en/article?id=zm_kb&sysparm_article=KB0068635}{sharing a slide-show window}.
The lecture PDFs themselves do not depend on Zoom.
\end{document}
